\documentclass[../../main.tex]{subfiles}

\begin{document}

\section{连续函数}\label{section:ch02-连续函数}

连续函数的概念是许多数学学科的基础，每一本有关微积分的书，都是首先讲述实直线上的连续函数，然后会提到平面和空间上的连续函数。而更一般的连续函数通常在数学的一些后继课程中讲述。本节将给出连续函数的定义（以上述各种连续函数为特例），我们还要研究连续函数的一些性质，其中的大部分内容都是我们在微积分和分析中所给出的连续函数的性质的直接推广。

\subsection*{函数的连续性}

\begin{definition}[连续函数]\label{definition:ch02-连续函数}
设 $X$ 和 $Y$ 是两个拓扑空间。函数 $f: X \rightarrow Y$ 称为连续的（continuous），如果对于 $Y$ 中的每一个开子集 $V$，$f^{-1}(V)$ 是 $X$ 中的一个开子集。
\end{definition}

注意，$f^{-1}(V)$ 是 $X$ 中所有使得 $f(x) \in V$ 的点 $x$ 的集合。如果 $V$ 与 $f$ 的像集 $f(X)$ 无交，那么 $f^{-1}(V) = \varnothing$。

函数的连续性不仅依赖于 $f$ 本身，也依赖于它的定义域和值域确定的拓扑。要想强调指出这一点，就说 $f$ 相对于 $X$ 的某拓扑和 $Y$ 的某拓扑而言是连续的。

注意，如果值域 $Y$ 的拓扑是由基 $\mathcal{B}$ 给出的，那么证明 $f$ 连续，就只要证明每一个基元素的原像是开的就行了。这是因为 $Y$ 的任意开集 $V$，可以写成基元素的并，即
\begin{align*}
V = \bigcup_{\alpha \in J} B_{\alpha}.
\end{align*}
因此
\begin{align*}
f^{-1}(V) = \bigcup_{\alpha \in J} f^{-1}(B_{\alpha}),
\end{align*}
如果每一个 $f^{-1}(B_{\alpha})$ 是开的，那么 $f^{-1}(V)$ 就是开的。

如果 $Y$ 的拓扑是由子基 $\mathcal{S}$ 给出的，那么为了证明 $f$ 的连续性，只要证明每一个子基元素的原像是开的就可以了。这是因为 $Y$ 的任意基元素 $B$ 可以写成子基元素的有限交 $S_{1} \cap \cdots \cap S_{n}$。由于
\begin{align*}
f^{-1}(B) = f^{-1}(S_{1}) \cap \cdots \cap f^{-1}(S_{n}),
\end{align*}
因此每一个基元素的原像是开的。

\begin{example}\label{example:ch02-eps-delta}
我们来看看在分析中研究过的一个函数，即实变量的实值函数
\begin{align*}
f: \mathbb{R} \longrightarrow \mathbb{R}.
\end{align*}
在分析中，采用“$\varepsilon$-$\delta$ 定义”来定义 $f$ 的连续性。多少年来，每一个初学数学的学生都对此感到困惑不解。其实，$\varepsilon$-$\delta$ 定义与上面的定义是等价的。例如，我们可以这样证明上面的定义蕴涵 $\varepsilon$-$\delta$ 定义：

对于 $\mathbb{R}$ 中的点 $x_{0}$，给定 $\varepsilon > 0$，区间 $V = (f(x_{0}) - \varepsilon, f(x_{0}) + \varepsilon)$ 是值域空间 $\mathbb{R}$ 的开集，因此 $f^{-1}(V)$ 是定义域空间 $\mathbb{R}$ 的开集。因为 $f^{-1}(V)$ 包含 $x_{0}$，所以必定包含 $x_{0}$ 的某一个基元素 $(a, b)$。选取 $\delta$ 为 $x_{0} - a$ 与 $b - x_{0}$ 两者中的较小者。那么，若 $|x - x_{0}| < \delta$，则点 $x$ 必定在 $(a, b)$ 之中，于是 $f(x) \in V$，从而得到 $|f(x) - f(x_{0})| < \varepsilon$。

要想证明 $\varepsilon$-$\delta$ 定义蕴涵我们上面的定义也不困难，我们留给读者去完成。当我们研究度量空间的时候，还要回过头来考虑这个例子。
\end{example}

\begin{example}\label{example:ch02-微积分函数}
在微积分中要研究很多种函数的连续性质，例如，我们研究下面几种类型的函数：
\begin{align*}
f: \mathbb{R} \longrightarrow \mathbb{R}^{2} \quad \text{（平面曲线）},
\end{align*}
\begin{align*}
f: \mathbb{R} \longrightarrow \mathbb{R}^{3} \quad \text{（空间曲线）},
\end{align*}
\begin{align*}
f: \mathbb{R}^{2} \longrightarrow \mathbb{R} \quad \text{（两个实变量的函数 } f(x, y)\text{）},
\end{align*}
\begin{align*}
f: \mathbb{R}^{3} \longrightarrow \mathbb{R} \quad \text{（三个实变量的函数 } f(x, y, z)\text{）},
\end{align*}
\begin{align*}
f: \mathbb{R}^{2} \longrightarrow \mathbb{R}^{2} \quad \text{（平面上的向量场 } \bm{v}(x, y)\text{）}.
\end{align*}
对于其中每一种函数都可以定义连续性。我们所给出的连续性的一般定义，包括了所有这些特殊情形，这不过是将要证明的有关积空间及度量空间上连续函数的一些一般性定理的一个推论。
\end{example}

\begin{example}\label{example:ch02-恒等映射}
设 $\mathbb{R}$ 为具有通常拓扑的实数集。$\mathbb{R}_{\ell}$ 为具有下限拓扑的实数集。令
\begin{align*}
f: \mathbb{R} \longrightarrow \mathbb{R}_{\ell}
\end{align*}
为恒等函数，即对于每一个实数 $x$，$f(x) = x$。但是 $f$ 不是一个连续函数。因为 $\mathbb{R}_{\ell}$ 的开集 $[a, b)$ 的原像还是 $[a, b)$，它不是 $\mathbb{R}$ 的开集。而另一方面，恒等函数
\begin{align*}
g: \mathbb{R}_{\ell} \longrightarrow \mathbb{R}
\end{align*}
却是连续的。因为 $(a, b)$ 的原像是 $(a, b)$，它也是 $\mathbb{R}_{\ell}$ 中的开集。
\end{example}

在分析中研究过连续函数的许多不同但又等价的定义。它们中的一些可以推广到任意空间上去，下面的定理就研究这个问题。但是，大家熟悉的“$\varepsilon$-$\delta$ 定义”及“收敛序列定义”却不能推广到任意空间。研究度量空间时，将要处理这些问题。

\begin{theorem}[连续函数的等价刻画]\label{theorem:ch02-连续等价}
设 $X$ 和 $Y$ 是两个拓扑空间，$f: X \rightarrow Y$。下列条件是等价的：
\begin{enumerate}[(1)]
\item $f$ 连续.
\item 对于 $X$ 的任意一个子集 $A$，有 $f(\overline{A}) \subset \overline{f(A)}$.
\item 对于 $Y$ 的任意一个闭集 $B$，$f^{-1}(B)$ 是 $X$ 中的一个闭集.
\item 对于每一个 $x \in X$ 和 $f(x)$ 的每一个邻域 $V$，存在 $x$ 的一个邻域 $U$ 使得 $f(U) \subset V$.
\end{enumerate}
对于 $X$ 中的点 $x$，如果条件 (4) 成立，则称 $f$ 在点 $x$ 连续。
\end{theorem}

\begin{proof}
我们将证明 (1) $\Rightarrow$ (2) $\Rightarrow$ (3) $\Rightarrow$ (1) 及 (1) $\Rightarrow$ (4) $\Rightarrow$ (1)。

(1) $\Rightarrow$ (2)。设 $f$ 连续。$A$ 是 $X$ 的一个子集。下面证明：若 $x \in \overline{A}$，则 $f(x) \in \overline{f(A)}$。设 $V$ 是 $f(x)$ 的一个邻域，则 $f^{-1}(V)$ 是 $X$ 中包含 $x$ 的一个开集，它必定与 $A$ 相交于某点 $y$，于是 $V$ 与 $f(A)$ 有交点 $f(y)$。因此得到 $f(x) \in \overline{f(A)}$。

(2) $\Rightarrow$ (3)。设 $B$ 是 $Y$ 的一个闭集，$A = f^{-1}(B)$。下面证明 $A$ 是 $X$ 的一个闭集，也就是证明 $\overline{A} \subset A$。由初等集合论可见 $f(A) \subset B$。因此，对于 $x \in \overline{A}$ 有
\begin{align*}
f(x) \in f(\overline{A}) \subset \overline{f(A)} \subset \overline{B} = B.
\end{align*}
于是 $x \in f^{-1}(B) = A$。所以得到 $\overline{A} \subset A$。

(3) $\Rightarrow$ (1)。设 $V$ 为 $Y$ 的一个开集，$B = Y - V$。那么
\begin{align*}
f^{-1}(B) = f^{-1}(Y) - f^{-1}(V) = X - f^{-1}(V).
\end{align*}
易见 $B$ 是 $Y$ 中的一个闭集。根据假设，$f^{-1}(B)$ 是 $X$ 的一个闭集。所以 $f^{-1}(V)$ 是开的。

(1) $\Rightarrow$ (4)。设 $x \in X$，$V$ 是 $f(x)$ 的一个邻域。则 $U = f^{-1}(V)$ 是 $x$ 的一个邻域，满足 $f(U) \subset V$。

(4) $\Rightarrow$ (1)。设 $V$ 是 $Y$ 的一个开集，$x$ 是 $f^{-1}(V)$ 的一个点。则 $f(x) \in V$，根据假设可见，存在 $x$ 的一个邻域 $U_{x}$ 使得 $f(U_{x}) \subset V$。从而 $U_{x} \subset f^{-1}(V)$。由此得到 $f^{-1}(V)$ 可以表示成所有这些开集 $U_{x}$ 之并，从而它是开的。
\end{proof}

\subsection*{同胚}

设 $X$ 和 $Y$ 都是拓扑空间，$f: X \rightarrow Y$ 是一个一一映射。如果函数 $f$ 和它的反函数
\begin{align*}
f^{-1}: Y \longrightarrow X
\end{align*}
都连续，则称 $f$ 为一个同胚（homeomorphism）。

\par
\begin{figure}[htbp]
\centering
\includegraphics[width=0.55\textwidth]{ch02-fig18-1}
\caption{同胚}\label{figure:ch02-同胚}
\end{figure}

$f^{-1}$ 连续这个条件是说，对于 $X$ 的每一个开集 $U$，在 $f^{-1}: Y \rightarrow X$ 下它的原像是 $Y$ 中的开集。但是，$U$ 在映射 $f^{-1}$ 下的原像就是 $U$ 在映射 $f$ 下的像，参见\cref{figure:ch02-同胚}。于是，可以用另一种方式将同胚定义为一一映射 $f: X \rightarrow Y$，使得 $f(U)$ 是一个开集的充分必要条件是 $U$ 是一个开集。

这个附注说明，一个同胚 $f: X \rightarrow Y$，不仅仅给出了 $X$ 和 $Y$ 之间的一个一一映射，还给出了 $X$ 与 $Y$ 的开集族之间的一个一一映射。因此，完全藉助于 $X$ 的拓扑（即藉助于 $X$ 的开集）所得出的 $X$ 的任何一条性质，通过对应 $f$ 就可以推出空间 $Y$ 的相应性质。$X$ 的这种性质称为 $X$ 的一个拓扑性质（topological property）。

读者可能已在近世代数学过代数对象之间同构的概念，比如群与群之间，环与环之间。一个同构是一个保持代数结构的一一映射。拓扑学中与此类似的概念就是同胚，它是保持拓扑结构的一一映射。

现在设 $f: X \rightarrow Y$ 是一个连续的单射，$X$ 和 $Y$ 是两个拓扑空间。用 $Z$ 表示像集 $f(X)$，把它看成 $Y$ 的一个子空间。那么由限制 $f$ 的值域得到的函数 $f': X \rightarrow Z$ 就是一一映射。若 $f'$ 正好是 $X$ 与 $Z$ 之间的一个同胚，则称映射 $f: X \rightarrow Y$ 是一个拓扑嵌入（topological imbedding），或者称为从 $X$ 到 $Y$ 中的一个嵌入（imbedding）。

\par
\begin{figure}[htbp]
\centering
\includegraphics[width=0.3\textwidth]{ch02-fig18-2}
\caption{$f(x) = 3x + 1$}\label{figure:ch02-线性同胚}
\end{figure}

\begin{example}\label{example:ch02-线性同胚}
由 $f(x) = 3x + 1$ 所给出的函数 $f: \mathbb{R} \rightarrow \mathbb{R}$ 是一个同胚，参见\cref{figure:ch02-线性同胚}。


考虑用
\begin{align*}
g(y) = \frac{1}{3}(y - 1)
\end{align*}
定义的函数 $g: \mathbb{R} \rightarrow \mathbb{R}$。容易验证：对于所有实数 $x$ 和 $y$ 有 $f(g(y)) = y$ 和 $g(f(x)) = x$，因此 $f$ 是一个一一映射并且 $g = f^{-1}$。至于 $f$ 与 $g$ 的连续性则是微积分中熟知的结论。
\end{example}

\par
\begin{figure}[htbp]
\centering
\includegraphics[width=0.35\textwidth]{ch02-fig18-3}
\caption{$F(x) = \frac{x}{1 - x^{2}}$}\label{figure:ch02-区间同胚}
\end{figure}

\begin{example}\label{example:ch02-区间同胚}
用
\begin{align*}
F(x) = \frac{x}{1 - x^{2}}
\end{align*}
定义的函数 $F: (-1, 1) \rightarrow \mathbb{R}$ 是一个同胚，参见\cref{figure:ch02-区间同胚}。我们已经在第\cref{section:ch01-关系}节的例 9 中知道，$F$ 是一个保序的一一映射，它的反函数 $G$ 定义为
\begin{align*}
G(y) = \frac{2y}{1 + (1 + 4y^{2})^{1/2}}.
\end{align*}


可以用两种方法证明 $F$ 是一个同胚。第一个方法应用 $F$ 是一个保序的一一映射，从而把 $(-1, 1)$ 中序拓扑的一个基元素变成 $\mathbb{R}$ 中序拓扑的一个基元素，反之亦然。这样 $F$ 自然就是 $(-1, 1)$ 与 $\mathbb{R}$ 之间（相对于各自的序拓扑而言）的一个同胚。因为 $(-1, 1)$ 上的序拓扑就是通常的（子空间）拓扑，所以 $F$ 是 $(-1, 1)$ 与 $\mathbb{R}$ 的一个同胚。

证明 $F$ 是同胚的第二个方法是应用微积分中熟悉的结果，即代数函数和根式函数的连续性来证明 $F$ 和 $G$ 都是连续函数。
\end{example}

\par
\begin{figure}[htbp]
\centering
\includegraphics[width=0.5\textwidth]{ch02-fig18-4}
\caption{连续但不同胚的映射}\label{figure:ch02-圆周映射}
\end{figure}

\begin{example}\label{example:ch02-连续非同胚}
一一映射 $f: X \rightarrow Y$ 可以是连续的而不是同胚。\cref{example:ch02-恒等映射}中研究过的恒等映射 $g: \mathbb{R}_{\ell} \rightarrow \mathbb{R}$ 就是这样一个函数。还可以举出一个例子：设 $S^{1}$ 表示单位圆周（unit circle），
\begin{align*}
S^{1} = \{x \times y \mid x^{2} + y^{2} = 1\}.
\end{align*}
将其作为平面 $\mathbb{R}^{2}$ 的一个子空间。设映射
\begin{align*}
f: [0, 1) \longrightarrow S^{1}
\end{align*}
定义为 $f(t) = (\cos 2\pi t, \sin 2\pi t)$。我们应用三角函数的性质推得 $f$ 是一个连续的一一映射。但是 $f^{-1}$ 不连续。例如，定义域中的开集 $U = \left[0, \frac{1}{4}\right)$ 在 $f$ 下的像不是 $S^{1}$ 中的开集，这是由于不存在 $\mathbb{R}^{2}$ 中包含 $p = f(0)$ 的开集 $V$，使得 $V \cap S^{1} \subset f(U)$，参见\cref{figure:ch02-圆周映射}。

\end{example}

\begin{example}\label{example:ch02-非嵌入}
考虑函数
\begin{align*}
g: [0, 1) \rightarrow \mathbb{R}^{2},
\end{align*}
它是把上例中函数 $f$ 的值域扩大而得出的。映射 $g$ 是一个连续单射，但不是嵌入。
\end{example}

\subsection*{构造连续函数}

怎样构造从一个拓扑空间到另一个拓扑空间的连续函数呢？对于在分析中使用过的那些方法，有一些可以推广到任意拓扑空间，有一些则不行。我们首先研究某些对于一般拓扑空间成立的构造法，其余的留到后面再讲。

\begin{theorem}[构造连续函数的法则]\label{theorem:ch02-连续构造法则}
设 $X$，$Y$ 和 $Z$ 都是拓扑空间。
\begin{enumerate}[(a)]
\item (常值函数) 若 $f: X \rightarrow Y$ 将整个 $X$ 映成 $Y$ 的一个点 $y_{0}$，则 $f$ 连续.
\item (内射) 若 $A$ 为 $X$ 的一个子空间，则内射 $j: A \rightarrow X$ 连续.
\item (复合) 若 $f: X \rightarrow Y$ 与 $g: Y \rightarrow Z$ 连续，则映射 $g \circ f: X \rightarrow Z$ 连续.
\item (限制定义域) 设 $f: X \rightarrow Y$ 连续，$A$ 为 $X$ 的一个子空间，则限制映射 $f \mid A: A \rightarrow Y$ 连续.
\item (限制或扩大值域) 设 $f: X \rightarrow Y$ 连续，$Z$ 为 $Y$ 中包含像集 $f(X)$ 的一个子空间，则限制 $f$ 的值域而得到的函数 $g: X \rightarrow Z$ 也连续。若 $Z$ 以 $Y$ 为其子空间，则扩大 $f$ 的值域而得到的函数 $h: X \rightarrow Z$ 也连续.
\item (连续性的局部表示) 如果 $X$ 可以写成开集 $U_{\alpha}$ 的并，使得对于每一个 $\alpha$，$f \mid U_{\alpha}$ 连续，则映射 $f: X \rightarrow Y$ 连续.
\end{enumerate}
\end{theorem}

\begin{proof}
(a) 对于 $X$ 的任意一个点 $x$，设 $f(x) = y_{0}$。设 $V$ 是 $Y$ 中的一个开集。若 $V$ 包含点 $y_{0}$，则 $f^{-1}(V) = X$。若 $V$ 不包含 $y_{0}$，则 $f^{-1}(V) = \varnothing$。无论哪一种情形都得到 $f^{-1}(V)$ 是开的。

(b) 若 $U$ 是 $X$ 的一个开集，则根据空间拓扑的定义，$j^{-1}(U) = U \cap A$ 是 $A$ 的一个开集。

(c) 若 $U$ 是 $Z$ 的一个开集，则 $g^{-1}(U)$ 是 $Y$ 的一个开集，$f^{-1}(g^{-1}(U))$ 是 $X$ 的一个开集，而根据初等集合论可见
\begin{align*}
f^{-1}(g^{-1}(U)) = (g \circ f)^{-1}(U).
\end{align*}

(d) 函数 $f \mid A$ 是内射 $j: A \rightarrow X$ 与映射 $f: X \rightarrow Y$ 的复合，而这两者都是连续函数。

(e) 设 $f: X \rightarrow Y$ 连续。若 $f(x) \subset Z \subset Y$，我们来证明由 $f$ 得到的函数 $g: X \rightarrow Z$ 连续。设 $B$ 是 $Z$ 的一个开集，则存在 $Y$ 的某一个开集 $U$，使得 $B = Z \cap U$。因为 $Z$ 包含整个像集 $f(X)$，根据初等集合论可见
\begin{align*}
f^{-1}(U) = g^{-1}(B).
\end{align*}
因为 $f^{-1}(U)$ 是开的，所以 $g^{-1}(B)$ 也是开的。

如果 $Z$ 以 $Y$ 为子空间，我们证明 $h: X \rightarrow Z$ 连续。而这只要注意到 $h$ 是映射 $f: X \rightarrow Y$ 与内射 $j: Y \rightarrow Z$ 的复合就行了。

(f) 根据假设，$X$ 可以写成开集 $U_{\alpha}$ 的并，并且对于每一个 $\alpha$，$f \mid U_{\alpha}$ 连续。设 $V$ 是 $Y$ 的一个开集，则
\begin{align*}
f^{-1}(V) \cap U_{\alpha} = (f \mid U_{\alpha})^{-1}(V).
\end{align*}
这是因为上式两边表示的都是 $U_{\alpha}$ 中满足 $f(x) \in V$ 的点 $x$ 的集合。因为 $f \mid U_{\alpha}$ 连续，所以上述集合是 $U_{\alpha}$ 中的开集，因此是 $X$ 中的开集。而
\begin{align*}
f^{-1}(V) = \bigcup_{\alpha} (f^{-1}(V) \cap U_{\alpha}).
\end{align*}
于是，$f^{-1}(V)$ 也是 $X$ 的一个开集。
\end{proof}

\begin{theorem}[黏结引理]\label{theorem:ch02-黏结引理}
设 $X = A \cup B$ 并且 $A$ 和 $B$ 都是 $X$ 中闭集。$f: A \rightarrow Y$ 与 $g: B \rightarrow Y$ 都是连续函数。若对于任意 $x \in A \cap B$ 有 $f(x) = g(x)$，则 $f$ 和 $g$ 可以组成一个连续函数 $h: X \rightarrow Y$，它定义为：当 $x \in A$ 时，$h(x) = f(x)$；当 $x \in B$ 时，$h(x) = g(x)$。
\end{theorem}

\begin{proof}
设 $C$ 为 $Y$ 的一个闭集，根据初等集合论有
\begin{align*}
h^{-1}(C) = f^{-1}(C) \cup g^{-1}(C).
\end{align*}
因为 $f$ 连续，所以 $f^{-1}(C)$ 是 $A$ 的一个闭集。类似地，$g^{-1}(C)$ 是 $B$ 的一个闭集，从而也是 $X$ 的一个闭集。于是它们的并 $h^{-1}(C)$ 是 $X$ 的一个闭集。
\end{proof}

若 $A$ 和 $B$ 都是 $X$ 的开集，这个定理仍然成立。这正是“连续性的局部表示”法则（见\cref{theorem:ch02-连续构造法则}）的一个特例。

\par
\begin{figure}[htbp]
\centering
\includegraphics[width=0.65\textwidth]{ch02-fig18-5}
\caption{黏结引理的例子}\label{figure:ch02-黏结函数}
\end{figure}

\begin{example}\label{example:ch02-黏结函数}
函数 $h: \mathbb{R} \rightarrow \mathbb{R}$ 定义为
\begin{align*}
h(x) =
\begin{cases}
x & \text{对于 } x \leqslant 0, \\
x/2 & \text{对于 } x \geqslant 0.
\end{cases}
\end{align*}
上式中 $h$ 在每一个“区间段”都是连续函数，并且它们在定义域的公共部分，即单点集 $\{0\}$ 上的函数值相同。由于两者的定义域都是 $\mathbb{R}$ 的闭集，所以 $h$ 连续。为了能定义这种函数，就要求对于各个“区间段”来说，函数在其定义域的重合处相等。例如下式
\begin{align*}
k(x) =
\begin{cases}
x - 2 & \text{对于 } x \leqslant 0, \\
x + 2 & \text{对于 } x \geqslant 0,
\end{cases}
\end{align*}
就不能定义一个函数。另一方面，对于集合 $A$ 和 $B$ 也要进行某些限制以保证函数的连续性，例如下式
\begin{align*}
l(x) =
\begin{cases}
x - 2 & \text{对于 } x < 0, \\
x + 2 & \text{对于 } x \geqslant 0,
\end{cases}
\end{align*}
定义了从 $\mathbb{R}$ 到 $\mathbb{R}$ 的一个函数，它的两个部分虽然是连续的，但是 $l$ 却不是连续函数。因为开集 $(1, 3)$ 的原像 $[0, 1)$ 不是开集。参见\cref{figure:ch02-黏结函数}。

\end{example}

\begin{theorem}[到积空间的映射]\label{theorem:ch02-积空间映射}
设 $f: A \rightarrow X \times Y$ 定义为
\begin{align*}
f(a) = (f_{1}(a), f_{2}(a)).
\end{align*}
则 $f$ 连续的充分必要条件是函数
\begin{align*}
f_{1}: A \longrightarrow X \quad \text{与} \quad f_{2}: A \longrightarrow Y
\end{align*}
都连续。

映射 $f_{1}$ 和 $f_{2}$ 称为 $f$ 的坐标函数（coordinate function）。
\end{theorem}

\begin{proof}
设 $\pi_{1}: X \times Y \rightarrow X$ 与 $\pi_{2}: X \times Y \rightarrow Y$ 分别是到第一个和第二个坐标空间上的投射。它们是连续的。这是因为若设 $U$ 和 $V$ 分别是 $X$ 和 $Y$ 中的开集，则 $\pi_{1}^{-1}(U) = U \times Y$ 和 $\pi_{2}^{-1}(V) = X \times V$ 都是开集。注意，对于每一个 $a \in A$，
\begin{align*}
f_{1}(a) = \pi_{1}(f(a)), \quad f_{2}(a) = \pi_{2}(f(a)).
\end{align*}

若 $f$ 是连续函数，则 $f_{1}$ 和 $f_{2}$ 是连续函数的复合，因而都是连续的。反之，设 $f_{1}$ 和 $f_{2}$ 连续，我们证明：对于 $X \times Y$ 的拓扑的每一个基元素 $U \times V$，其原像 $f^{-1}(U \times V)$ 是开集。点 $a$ 在 $f^{-1}(U \times V)$ 中当且仅当 $f(a) \in U \times V$，也就是当且仅当 $f_{1}(a) \in U$ 并且 $f_{2}(a) \in V$。因此
\begin{align*}
f^{-1}(U \times V) = f_{1}^{-1}(U) \cap f_{2}^{-1}(V).
\end{align*}
由于 $f_{1}^{-1}(U)$ 和 $f_{2}^{-1}(V)$ 都是开集，所以它们的交也是开的。
\end{proof}

对于定义域是积空间的映射 $f: A \times B \rightarrow X$，没有一个常用的方法来判断其连续性。有这样一种猜想：如果 $f$“分别关于每一个变量”连续，则 $f$ 连续。但它是不对的（见习题 12）。

\begin{example}\label{example:ch02-参数曲线}
在微积分中，平面上一条参数曲线被定义为一个连续映射 $f: [a, b] \rightarrow \mathbb{R}^{2}$，通常写成 $f(t) = (x(t), y(t))$。我们常常要用到这样的结论：若 $x$ 和 $y$ 都是 $t$ 的连续函数，则 $f$ 是 $t$ 的一个连续函数。类似地，平面向量场
\begin{align*}
\bm{v}(x, y) = P(x, y)\bm{i} + Q(x, y)\bm{j} = (P(x, y), Q(x, y))
\end{align*}
称为连续的，如果 $P$ 和 $Q$ 都是连续函数；或者等价地说，$\bm{v}$ 作为 $\mathbb{R}^{2}$ 到 $\mathbb{R}^{2}$ 的映射是连续的。这两种说法都是上面定理的特殊情形。
\end{example}

在分析中大量使用的构造连续函数的方法是取连续实值函数的和、差、积、商。有一个标准定理：若 $f, g: X \rightarrow \mathbb{R}$ 连续，则 $f + g$、$f - g$、$f \cdot g$ 都连续，并且对于所有使 $g(x) \neq 0$ 的 $x$，$f/g$ 连续。我们将在第\cref{section:ch02-度量拓扑续}节研究这个定理。

在分析中我们所熟悉的构造连续函数的另一个办法，是取函数的无穷序列的极限。有这样一个定理，其大意是说，如果一个实变量连续实值函数的序列一致收敛于一个极限函数，那么极限函数必为连续函数。这个定理称为一致极限定理。例如，当正弦、余弦函数的定义是严格地用无穷级数给出的时候，用这个定理可以证明三角函数的连续性。这个定理能够推广为关于从任意拓扑空间 $X$ 到度量空间 $Y$ 中映射的定理，我们将在第\cref{section:ch02-度量拓扑续}节予以证明。

\subsection*{习题}

\begin{exercise}\label{exercise:ch02-连续函数-1}
证明：函数 $f: \mathbb{R} \rightarrow \mathbb{R}$ 连续性的 $\varepsilon$-$\delta$ 定义蕴涵开集定义。
\end{exercise}

\begin{exercise}\label{exercise:ch02-连续函数-2}
设 $f: X \rightarrow Y$ 连续，$x$ 是 $X$ 的子集 $A$ 的一个极限点，那么 $f(x)$ 一定是 $f(A)$ 的极限点吗？
\end{exercise}

\begin{exercise}\label{exercise:ch02-连续函数-3}
设 $X$ 和 $X'$ 分别表示具有拓扑 $\mathcal{T}$ 和 $\mathcal{T}'$ 的同一个集合，$i: X' \rightarrow X$ 为恒等函数。

(a) 证明：$i$ 连续 $\Longleftrightarrow \mathcal{T}'$ 细于 $\mathcal{T}$。

(b) 证明：$i$ 是同胚 $\Longleftrightarrow \mathcal{T}' = \mathcal{T}$。
\end{exercise}

\begin{exercise}\label{exercise:ch02-连续函数-4}
给定 $x_{0} \in X$ 和 $y_{0} \in Y$，证明：分别由
\begin{align*}
f(x) = x \times y_{0} \quad \text{和} \quad g(y) = x_{0} \times y
\end{align*}
所定义的映射 $f: X \rightarrow X \times Y$ 与 $g: Y \rightarrow X \times Y$ 都是嵌入。
\end{exercise}

\begin{exercise}\label{exercise:ch02-连续函数-5}
证明：$\mathbb{R}$ 的子空间 $(a, b)$ 与 $(0, 1)$ 同胚；$\mathbb{R}$ 的子空间 $[a, b]$ 与 $[0, 1]$ 同胚。
\end{exercise}

\begin{exercise}\label{exercise:ch02-连续函数-6}
找出一个只在一点连续的函数 $f: \mathbb{R} \rightarrow \mathbb{R}$。
\end{exercise}

\begin{exercise}\label{exercise:ch02-连续函数-7}
(a) 假定 $f: \mathbb{R} \rightarrow \mathbb{R}$“右连续”，即对于每一个 $a \in \mathbb{R}$，
\begin{align*}
\lim_{x \rightarrow a^{+}} f(x) = f(a),
\end{align*}
证明：将 $f$ 看成从 $\mathbb{R}_{\ell}$ 到 $\mathbb{R}$ 的一个函数时它是连续的。

(b) 如果把函数 $f: \mathbb{R} \rightarrow \mathbb{R}$ 看成从 $\mathbb{R}$ 到 $\mathbb{R}_{\ell}$ 的映射时，什么样的 $f$ 是连续的？如果看成从 $\mathbb{R}_{\ell}$ 到 $\mathbb{R}_{\ell}$ 的映射呢？我们将在第 3 章中讨论这个问题。
\end{exercise}

\begin{exercise}\label{exercise:ch02-连续函数-8}
设 $Y$ 为具有序拓扑的全序集，$f, g: X \rightarrow Y$ 是连续的。

(a) 证明：$\{x \mid f(x) \leqslant g(x)\}$ 在 $X$ 中是闭的。

(b) 证明：由下式
\begin{align*}
h(x) = \min\{f(x), g(x)\}
\end{align*}
所定义的函数 $h: X \rightarrow Y$ 是连续的。[提示：应用黏结引理。]
\end{exercise}

\begin{exercise}\label{exercise:ch02-连续函数-9}
设 $\{A_{\alpha}\}$ 是 $X$ 的一个子集族，$X = \bigcup_{\alpha} A_{\alpha}$。设 $f: X \rightarrow Y$，对于每一个 $\alpha$，$f \mid A_{\alpha}$ 连续。

(a) 若 $\{A_{\alpha}\}$ 为有限族，并且每一个 $A_{\alpha}$ 都是闭集，则 $f$ 连续。

(b) 找出一个可数族 $\{A_{\alpha}\}$，其中每一个 $A_{\alpha}$ 都是闭集，但 $f$ 不连续的例子。

(c) 一个加标集族 $\{A_{\alpha}\}$ 称为局部有限的（locally finite），如果 $X$ 的每一个点 $x$ 都有一个邻域，仅与有限多个 $A_{\alpha}$ 相交。证明：若族 $\{A_{\alpha}\}$ 是局部有限的，并且每一个 $A_{\alpha}$ 都是闭集，则 $f$ 连续。
\end{exercise}

\begin{exercise}\label{exercise:ch02-连续函数-10}
设 $f: A \rightarrow B$ 和 $g: C \rightarrow D$ 都是连续函数。证明：由公式
\begin{align*}
(f \times g)(a \times c) = f(a) \times g(c)
\end{align*}
所定义的映射 $f \times g: A \times C \rightarrow B \times D$ 是连续的。
\end{exercise}

\begin{exercise}\label{exercise:ch02-连续函数-11}
设 $F: X \times Y \rightarrow Z$。我们把 $F$ 分别关于每一个变量连续（continuous in each variable separately）定义为：对于 $Y$ 中每一个 $y_{0}$，用 $h(x) = F(x \times y_{0})$ 定义的映射 $h: X \rightarrow Z$ 连续，并且对于 $X$ 中每一个 $x_{0}$，用 $k(y) = F(x_{0} \times y)$ 所定义的映射 $k: Y \rightarrow Z$ 连续。证明：若 $F$ 连续，则 $F$ 分别关于每一个变量连续。
\end{exercise}

\begin{exercise}\label{exercise:ch02-连续函数-12}
设 $F: \mathbb{R} \times \mathbb{R} \rightarrow \mathbb{R}$ 定义为：
\begin{align*}
F(x \times y) =
\begin{cases}
xy/(x^{2} + y^{2}) & \text{如果 } x \times y \neq 0 \times 0, \\
0 & \text{如果 } x \times y = 0 \times 0.
\end{cases}
\end{align*}
(a) 证明：$F$ 分别关于每一个变量连续。

(b) 计算出由 $g(x) = F(x \times x)$ 所定义的函数 $g: \mathbb{R} \rightarrow \mathbb{R}$。

(c) 证明 $F$ 不连续。
\end{exercise}

\begin{exercise}\label{exercise:ch02-连续函数-13}
设 $A \subset X$，$f: A \rightarrow Y$ 连续，$Y$ 是一个 Hausdorff 空间。证明：若 $f$ 能扩充为一个连续函数 $g: \overline{A} \rightarrow Y$，则 $g$ 由 $f$ 唯一决定。
\end{exercise}

\end{document}